\IfFileExists{stacks-project.cls}{%
\documentclass{stacks-project}
}{%
\documentclass{amsart}
}

% For dealing with references we use the comment environment
\usepackage{verbatim}
\newenvironment{reference}{\comment}{\endcomment}
%\newenvironment{reference}{}{}
\newenvironment{slogan}{\comment}{\endcomment}
\newenvironment{history}{\comment}{\endcomment}

% For commutative diagrams we use Xy-pic
\usepackage[all]{xy}

% We use 2cell for 2-commutative diagrams.
\xyoption{2cell}
\UseAllTwocells

% We use multicol for the list of chapters between chapters
\usepackage{multicol}

% This is generally recommended for better output
\usepackage{lmodern}
\usepackage[T1]{fontenc}

% For cross-file-references
\usepackage{xr-hyper}

% Package for hypertext links:
\usepackage{hyperref}

% For any local file, say "hello.tex" you want to link to please
% use \externaldocument[hello-]{hello}
\externaldocument[introduction-]{introduction}
\externaldocument[conventions-]{conventions}
\externaldocument[sets-]{sets}
\externaldocument[categories-]{categories}
\externaldocument[topology-]{topology}
\externaldocument[sheaves-]{sheaves}
\externaldocument[sites-]{sites}
\externaldocument[stacks-]{stacks}
\externaldocument[fields-]{fields}
\externaldocument[algebra-]{algebra}
\externaldocument[brauer-]{brauer}
\externaldocument[homology-]{homology}
\externaldocument[derived-]{derived}
\externaldocument[simplicial-]{simplicial}
\externaldocument[more-algebra-]{more-algebra}
\externaldocument[smoothing-]{smoothing}
\externaldocument[modules-]{modules}
\externaldocument[sites-modules-]{sites-modules}
\externaldocument[injectives-]{injectives}
\externaldocument[cohomology-]{cohomology}
\externaldocument[sites-cohomology-]{sites-cohomology}
\externaldocument[dga-]{dga}
\externaldocument[dpa-]{dpa}
\externaldocument[sdga-]{sdga}
\externaldocument[hypercovering-]{hypercovering}
\externaldocument[schemes-]{schemes}
\externaldocument[constructions-]{constructions}
\externaldocument[properties-]{properties}
\externaldocument[morphisms-]{morphisms}
\externaldocument[coherent-]{coherent}
\externaldocument[divisors-]{divisors}
\externaldocument[limits-]{limits}
\externaldocument[varieties-]{varieties}
\externaldocument[topologies-]{topologies}
\externaldocument[descent-]{descent}
\externaldocument[perfect-]{perfect}
\externaldocument[more-morphisms-]{more-morphisms}
\externaldocument[flat-]{flat}
\externaldocument[groupoids-]{groupoids}
\externaldocument[more-groupoids-]{more-groupoids}
\externaldocument[etale-]{etale}
\externaldocument[chow-]{chow}
\externaldocument[intersection-]{intersection}
\externaldocument[pic-]{pic}
\externaldocument[weil-]{weil}
\externaldocument[adequate-]{adequate}
\externaldocument[dualizing-]{dualizing}
\externaldocument[duality-]{duality}
\externaldocument[discriminant-]{discriminant}
\externaldocument[derham-]{derham}
\externaldocument[local-cohomology-]{local-cohomology}
\externaldocument[algebraization-]{algebraization}
\externaldocument[curves-]{curves}
\externaldocument[resolve-]{resolve}
\externaldocument[models-]{models}
\externaldocument[functors-]{functors}
\externaldocument[equiv-]{equiv}
\externaldocument[pione-]{pione}
\externaldocument[etale-cohomology-]{etale-cohomology}
\externaldocument[proetale-]{proetale}
\externaldocument[relative-cycles-]{relative-cycles}
\externaldocument[more-etale-]{more-etale}
\externaldocument[trace-]{trace}
\externaldocument[crystalline-]{crystalline}
\externaldocument[spaces-]{spaces}
\externaldocument[spaces-properties-]{spaces-properties}
\externaldocument[spaces-morphisms-]{spaces-morphisms}
\externaldocument[decent-spaces-]{decent-spaces}
\externaldocument[spaces-cohomology-]{spaces-cohomology}
\externaldocument[spaces-limits-]{spaces-limits}
\externaldocument[spaces-divisors-]{spaces-divisors}
\externaldocument[spaces-over-fields-]{spaces-over-fields}
\externaldocument[spaces-topologies-]{spaces-topologies}
\externaldocument[spaces-descent-]{spaces-descent}
\externaldocument[spaces-perfect-]{spaces-perfect}
\externaldocument[spaces-more-morphisms-]{spaces-more-morphisms}
\externaldocument[spaces-flat-]{spaces-flat}
\externaldocument[spaces-groupoids-]{spaces-groupoids}
\externaldocument[spaces-more-groupoids-]{spaces-more-groupoids}
\externaldocument[bootstrap-]{bootstrap}
\externaldocument[spaces-pushouts-]{spaces-pushouts}
\externaldocument[spaces-chow-]{spaces-chow}
\externaldocument[groupoids-quotients-]{groupoids-quotients}
\externaldocument[spaces-more-cohomology-]{spaces-more-cohomology}
\externaldocument[spaces-simplicial-]{spaces-simplicial}
\externaldocument[spaces-duality-]{spaces-duality}
\externaldocument[formal-spaces-]{formal-spaces}
\externaldocument[restricted-]{restricted}
\externaldocument[spaces-resolve-]{spaces-resolve}
\externaldocument[formal-defos-]{formal-defos}
\externaldocument[defos-]{defos}
\externaldocument[cotangent-]{cotangent}
\externaldocument[examples-defos-]{examples-defos}
\externaldocument[algebraic-]{algebraic}
\externaldocument[examples-stacks-]{examples-stacks}
\externaldocument[stacks-sheaves-]{stacks-sheaves}
\externaldocument[criteria-]{criteria}
\externaldocument[artin-]{artin}
\externaldocument[quot-]{quot}
\externaldocument[stacks-properties-]{stacks-properties}
\externaldocument[stacks-morphisms-]{stacks-morphisms}
\externaldocument[stacks-limits-]{stacks-limits}
\externaldocument[stacks-cohomology-]{stacks-cohomology}
\externaldocument[stacks-perfect-]{stacks-perfect}
\externaldocument[stacks-introduction-]{stacks-introduction}
\externaldocument[stacks-more-morphisms-]{stacks-more-morphisms}
\externaldocument[stacks-geometry-]{stacks-geometry}
\externaldocument[moduli-]{moduli}
\externaldocument[moduli-curves-]{moduli-curves}
\externaldocument[examples-]{examples}
\externaldocument[exercises-]{exercises}
\externaldocument[guide-]{guide}
\externaldocument[desirables-]{desirables}
\externaldocument[coding-]{coding}
\externaldocument[obsolete-]{obsolete}
\externaldocument[fdl-]{fdl}
\externaldocument[index-]{index}

% Theorem environments.
%
\theoremstyle{plain}
\newtheorem{theorem}[subsection]{Theorem}
\newtheorem{proposition}[subsection]{Proposition}
\newtheorem{lemma}[subsection]{Lemma}

\theoremstyle{definition}
\newtheorem{definition}[subsection]{Definition}
\newtheorem{example}[subsection]{Example}
\newtheorem{exercise}[subsection]{Exercise}
\newtheorem{situation}[subsection]{Situation}

\theoremstyle{remark}
\newtheorem{remark}[subsection]{Remark}
\newtheorem{remarks}[subsection]{Remarks}

\numberwithin{equation}{subsection}

% Macros
%
\def\lim{\mathop{\mathrm{lim}}\nolimits}
\def\colim{\mathop{\mathrm{colim}}\nolimits}
\def\Spec{\mathop{\mathrm{Spec}}}
\def\Hom{\mathop{\mathrm{Hom}}\nolimits}
\def\Ext{\mathop{\mathrm{Ext}}\nolimits}
\def\SheafHom{\mathop{\mathcal{H}\!\mathit{om}}\nolimits}
\def\SheafExt{\mathop{\mathcal{E}\!\mathit{xt}}\nolimits}
\def\Sch{\mathit{Sch}}
\def\Mor{\mathop{\mathrm{Mor}}\nolimits}
\def\Ob{\mathop{\mathrm{Ob}}\nolimits}
\def\Sh{\mathop{\mathit{Sh}}\nolimits}
\def\NL{\mathop{N\!L}\nolimits}
\def\CH{\mathop{\mathrm{CH}}\nolimits}
\def\proetale{{pro\text{-}\acute{e}tale}}
\def\etale{{\acute{e}tale}}
\def\QCoh{\mathit{QCoh}}
\def\Ker{\mathop{\mathrm{Ker}}}
\def\Im{\mathop{\mathrm{Im}}}
\def\Coker{\mathop{\mathrm{Coker}}}
\def\Coim{\mathop{\mathrm{Coim}}}

% Boxtimes
%
\DeclareMathSymbol{\boxtimes}{\mathbin}{AMSa}{"02}

%
% Macros for moduli stacks/spaces
%
\def\QCohstack{\mathcal{QC}\!\mathit{oh}}
\def\Cohstack{\mathcal{C}\!\mathit{oh}}
\def\Spacesstack{\mathcal{S}\!\mathit{paces}}
\def\Quotfunctor{\mathrm{Quot}}
\def\Hilbfunctor{\mathrm{Hilb}}
\def\Curvesstack{\mathcal{C}\!\mathit{urves}}
\def\Polarizedstack{\mathcal{P}\!\mathit{olarized}}
\def\Complexesstack{\mathcal{C}\!\mathit{omplexes}}
% \Pic is the operator that assigns to X its picard group, usage \Pic(X)
% \Picardstack_{X/B} denotes the Picard stack of X over B
% \Picardfunctor_{X/B} denotes the Picard functor of X over B
\def\Pic{\mathop{\mathrm{Pic}}\nolimits}
\def\Picardstack{\mathcal{P}\!\mathit{ic}}
\def\Picardfunctor{\mathrm{Pic}}
\def\Deformationcategory{\mathcal{D}\!\mathit{ef}}
\usepackage{cleveref}
\crefname{originalreference}{original reference}{original references}
\Crefname{originalreference}{Original reference}{Original references}
\makeatletter
\newlabel{lemma-morphism-augmentation-simplicial-sites}{{0D99}{source}{}{sourcepointer-1}{}}
\newlabel{lemma-morphism-augmentation-simplicial-sites@cref}{{[originalreference][0][]0D99}{[1][0][]source}{}{}{}}
\makeatother
\crefname{originalreference}{original reference}{original references}
\Crefname{originalreference}{Original reference}{Original references}
\makeatletter
\newlabel{lemma-sr-when-fibre-products}{{09WL}{source}{}{sourcepointer-2}{}}
\newlabel{lemma-sr-when-fibre-products@cref}{{[originalreference][0][]09WL}{[1][0][]source}{}{}{}}
\makeatother
\crefname{originalreference}{original reference}{original references}
\Crefname{originalreference}{Original reference}{Original references}
\makeatletter
\newlabel{remark-morphism-augmentation-simplicial-sites}{{0D98}{source}{}{sourcepointer-3}{}}
\newlabel{remark-morphism-augmentation-simplicial-sites@cref}{{[originalreference][0][]0D98}{[1][0][]source}{}{}{}}
\makeatother
\crefname{originalsection}{original section}{original sections}
\Crefname{originalsection}{Original section}{Original sections}
\makeatletter
\newlabel{section-fppf-hypercovering}{{0DH4}{source}{}{sourcepointer-4}{}}
\newlabel{section-fppf-hypercovering@cref}{{[originalsection][0][]0DH4}{[1][0][]source}{}{}{}}
\makeatother
\crefname{originalsection}{original section}{original sections}
\Crefname{originalsection}{Original section}{Original sections}
\makeatletter
\newlabel{section-hypercovering}{{09A0}{source}{}{sourcepointer-5}{}}
\newlabel{section-hypercovering@cref}{{[originalsection][0][]09A0}{[1][0][]source}{}{}{}}
\makeatother
\crefname{originalsection}{original section}{original sections}
\Crefname{originalsection}{Original section}{Original sections}
\makeatletter
\newlabel{section-simplicial-algebraic-spaces}{{0DE7}{source}{}{sourcepointer-6}{}}
\newlabel{section-simplicial-algebraic-spaces@cref}{{[originalsection][0][]0DE7}{[1][0][]source}{}{}{}}
\makeatother
\crefname{originalreference}{original reference}{original references}
\Crefname{originalreference}{Original reference}{Original references}
\makeatletter
\newlabel{situation-simplicial-site}{{09WE}{source}{}{sourcepointer-7}{}}
\newlabel{situation-simplicial-site@cref}{{[originalreference][0][]09WE}{[1][0][]source}{}{}{}}
\makeatother
\crefname{originalsection}{original section}{original sections}
\Crefname{originalsection}{Original section}{Original sections}
\makeatletter
\newlabel{spaces-more-cohomology-section-ph-etale}{{0DGS}{source}{}{sourcepointer-8}{}}
\newlabel{spaces-more-cohomology-section-ph-etale@cref}{{[originalsection][0][]0DGS}{[1][0][]source}{}{}{}}
\makeatother
\title{Comparison diagram of topoi between the big ph simplicial site and the small etale site}
\author{The Stacks Project Authors}
\date{stacks-project:spaces-simplicial.tex:local-snapshot-20261012; GNU FDL 1.2 or later}
\begin{document}
\maketitle
\noindent\textit{Editorial note.} All delivered mathematical text is the author's own, copied from the pinned Stacks Project chapter file sources/stacks-spaces-simplicial/pinned-original-source.tex; no statement or proof step was rewritten, reordered or removed.\par\medskip
\noindent\textit{Editorial note.} The counted proof is the author's complete argument: the case-A assertions ('This is immediate from the construction', source lines 8007-8008) refer to the definitional form of the sites' coverings; the author's footnote distinguishing the ph total site from the fppf total site is preserved verbatim; this lemma underlies the descent theorems delivered as cards L46151 (descent of sheaves for ph hypercoverings) and L46163 (descent of sheaves for proper hypercoverings), and is the ph counterpart of card L46173 (fppf-etale comparison).\par\medskip
\noindent\textit{Editorial note.} The label lemma-compare-simplicial-objects-ph-etale names the counted statement, and the proof environment immediately following it in the source is delivered verbatim as its proof.\par\medskip
\noindent\textit{Editorial note.} The counted statement is carried by the author's own lemma environment in the chapter source.\par\medskip
\noindent\textit{Allowed standard prerequisites (editorial).}\begin{itemize}
\item The hypercovering total-site construction over (Spaces/S)\_ph. Same chapter, Section \textbackslash{}label\{section-hypercovering\}. Proof lines 7992-7998: the notation (Spaces/U)\_\{ph,total\} means the total-site construction applied to the site (Spaces/S)\_ph and the simplicial object U (footnote distinguishing it from the fppf total site).
\item The sites representing the small etale topoi. Same chapter, Section \textbackslash{}label\{section-simplicial-algebraic-spaces\}. Proof lines 7999-8001: X\_\{spaces,etale\} and U\_\{spaces,etale\} may be used for the right-hand topoi.
\item Both sites fall into case A of the simplicial-site situation. Same chapter, Situation \textbackslash{}label\{situation-simplicial-site\} and \textbackslash{}label\{lemma-sr-when-fibre-products\}. Proof lines 8004-8010: the case-A placement of (Spaces/U)\_\{ph,total\} and U\_\{spaces,etale\}.
\item The component morphisms of sites a\_\{U\_n\} and a\_X. More on Cohomology of Spaces chapter, Section \textbackslash{}label\{spaces-more-cohomology-section-ph-etale\}. Proof lines 8011-8019: U\_\{n,spaces,etale\} -> (Spaces/U\_n)\_ph and X\_\{spaces,etale\} -> (Spaces/X)\_ph are morphisms of sites, denoted a\_\{U\_n\} = epsilon\_\{U\_n\} pi\_\{u\_n\} and a\_X = epsilon\_X pi\_X.
\item Morphisms of simplicial sites compatible with augmentations. Same chapter, Remark \textbackslash{}label\{remark-morphism-augmentation-simplicial-sites\} and \textbackslash{}label\{lemma-morphism-augmentation-simplicial-sites\}. Proof lines 8020-8024: the componentwise data form a morphism of simplicial sites compatible with the augmentations, and the cited lemma produces the commutative diagram.
\end{itemize}\medskip
\ifdefined\hypertarget\hypertarget{sourcepointer-1}{}\fi
\noindent\textit{Source reference.} Original reference 0D99: sources/stacks-spaces-simplicial/pinned-original-source.tex, source lines 1402--1402. Prerequisite pointer invoked at proof line(s) 8023; it names a standard fact of the same chapter that the counted argument uses, and every step of the counted proof built on it is retained verbatim in the delivered body.\par\medskip
\ifdefined\hypertarget\hypertarget{sourcepointer-2}{}\fi
\noindent\textit{Source reference.} Original reference 09WL: sources/stacks-spaces-simplicial/pinned-original-source.tex, source lines 5136--5136. Prerequisite pointer invoked at proof line(s) 8009; it names a standard fact of the same chapter that the counted argument uses, and every step of the counted proof built on it is retained verbatim in the delivered body.\par\medskip
\ifdefined\hypertarget\hypertarget{sourcepointer-3}{}\fi
\noindent\textit{Source reference.} Original reference 0D98: sources/stacks-spaces-simplicial/pinned-original-source.tex, source lines 1372--1372. Prerequisite pointer invoked at proof line(s) 8021; it names a standard fact of the same chapter that the counted argument uses, and every step of the counted proof built on it is retained verbatim in the delivered body.\par\medskip
\ifdefined\hypertarget\hypertarget{sourcepointer-4}{}\fi
\noindent\textit{Source reference.} Original section 0DH4: sources/stacks-spaces-simplicial/pinned-original-source.tex, source lines 7188--7188. Prerequisite pointer invoked at proof line(s) 7998; it names a standard fact of the same chapter that the counted argument uses, and every step of the counted proof built on it is retained verbatim in the delivered body.\par\medskip
\ifdefined\hypertarget\hypertarget{sourcepointer-5}{}\fi
\noindent\textit{Source reference.} Original section 09A0: sources/stacks-spaces-simplicial/pinned-original-source.tex, source lines 5012--5012. Prerequisite pointer invoked at proof line(s) 7986; it names a standard fact of the same chapter that the counted argument uses, and every step of the counted proof built on it is retained verbatim in the delivered body.\par\medskip
\ifdefined\hypertarget\hypertarget{sourcepointer-6}{}\fi
\noindent\textit{Source reference.} Original section 0DE7: sources/stacks-spaces-simplicial/pinned-original-source.tex, source lines 6985--6985. Prerequisite pointer invoked at proof line(s) 7988; it names a standard fact of the same chapter that the counted argument uses, and every step of the counted proof built on it is retained verbatim in the delivered body.\par\medskip
\ifdefined\hypertarget\hypertarget{sourcepointer-7}{}\fi
\noindent\textit{Source reference.} Original reference 09WE: sources/stacks-spaces-simplicial/pinned-original-source.tex, source lines 888--888. Prerequisite pointer invoked at proof line(s) 8006; it names a standard fact of the same chapter that the counted argument uses, and every step of the counted proof built on it is retained verbatim in the delivered body.\par\medskip
\ifdefined\hypertarget\hypertarget{sourcepointer-8}{}\fi
\noindent\textit{Source reference.} Original section 0DGS: sources/stacks-chow/tags.txt, source lines 16882--16882. Prerequisite pointer invoked at proof line(s) 8017; it names a standard fact of a sibling chapter that the counted argument uses, and every step of the counted proof built on it is retained verbatim in the delivered body.\par\medskip
\begin{lemma}
\label{lemma-compare-simplicial-objects-ph-etale}
Let $S$ be a scheme. Let $X$ be an algebraic space over $S$.
Let $U$ be a simplicial algebraic space over $S$.
Let $a : U \to X$ be an augmentation. There is a commutative diagram
$$
\xymatrix{
\Sh((\textit{Spaces}/U)_{ph, total}) \ar[r]_-h \ar[d]_{a_{ph}} &
\Sh(U_\etale) \ar[d]^a \\
\Sh((\textit{Spaces}/X)_{ph}) \ar[r]^-{h_{-1}} &
\Sh(X_\etale)
}
$$
where the left vertical arrow is defined in
Section \ref{section-hypercovering}
and the right vertical arrow is defined in
Section \ref{section-simplicial-algebraic-spaces}.
\end{lemma}

\begin{proof}
The notation $(\textit{Spaces}/U)_{ph, total}$ indicates that
we are using the construction of
Section \ref{section-hypercovering}
for the site $(\textit{Spaces}/S)_{ph}$ and the
simplicial object $U$ of this site\footnote{To distinguish from
$(\textit{Spaces}/U)_{fppf, total}$ defined using the fppf
topology in Section \ref{section-fppf-hypercovering}.}.
We will use the sites $X_{spaces, \etale}$ and $U_{spaces, \etale}$
for the topoi on the right hand side; this is permissible
see discussion in Section \ref{section-simplicial-algebraic-spaces}.

\medskip\noindent
Observe that both $(\textit{Spaces}/U)_{ph, total}$ and
$U_{spaces, \etale}$
fall into case A of Situation \ref{situation-simplicial-site}.
This is immediate from the construction of
$U_\etale$ in Section \ref{section-simplicial-algebraic-spaces}
and it follows from Lemma \ref{lemma-sr-when-fibre-products}
for $(\textit{Spaces}/U)_{ph, total}$.
Next, consider the functors
$U_{n, spaces, \etale} \to (\textit{Spaces}/U_n)_{ph}$, $U \mapsto U/U_n$
and
$X_{spaces, \etale} \to (\textit{Spaces}/X)_{ph}$, $U \mapsto U/X$.
We have seen that these define morphisms of sites in
More on Cohomology of Spaces, Section
\ref{spaces-more-cohomology-section-ph-etale}
where these were denoted $a_{U_n} = \epsilon_{U_n} \circ \pi_{u_n}$
and $a_X = \epsilon_X \circ \pi_X$.
Thus we obtain a morphism of simplicial sites compatible with
augmentations as in Remark \ref{remark-morphism-augmentation-simplicial-sites}
and we may apply
Lemma \ref{lemma-morphism-augmentation-simplicial-sites} to conclude.
\end{proof}

\end{document}
